% Stage 8 restored fragment — 2.7.1 Null spaces of radiation operators (archive lines 9114-9365)
\subsection{Null spaces of radiation operators}
\label{sec:ch271}

On a declared compact-range or anechoic campaign, technicians record frequency, outline, and probe geometry before any harmonic fit is trusted \cite{hansen1988}. Non-radiating currents store reactive energy without far-zone export \cite{stratton1941}. Near-field richness must not be reported as far-zone diversity \cite{hansen1988}. Maxwell's retarded export to $\Gamma_R$ remains the governing constraint \cite{maxwell1873}.
Section~\ref{sec:ch27} opened with laboratory signatures---balanced loops, counter-directed
pairs, below-cutoff guides---and named the radiation kernel in plain language. Subsection
~\ref{sec:ch271} now supplies the operator account: structural properties of
$\ker\mathcal{R}_\omega$, adjoint null spaces, cokernel pairings, and modal null sectors,
developing Eq.~\eqref{eq:ch2.ker-R-omega} from Subsection~\ref{sec:ch244}
\cite{stratton1941}. The loop and scanner patterns cited in the section introduction are not
re-derived here; they motivate the locked definition below.

Assumptions. Fix \(\omega>0\), admissible source space \(\mathcal{J}_\Omega\)
Eq.~\eqref{eq:ch2.source-space-refined}, charge continuity and boundary compatibility
Eq.~\eqref{eq:ch2.source-admissibility}, outgoing radiation condition on exterior problems
Eqs.~\eqref{eq:ch2.sommerfeld-outgoing-class}--\eqref{eq:ch2.outgoing-radiation-refinement},
and radiation map \(\mathcal{R}_\omega=\Pi_{\mathrm{rad}}\circ\mathcal{S}_\omega\) from
Eqs.~\eqref{eq:ch2.radiation-operator-def}--\eqref{eq:ch2.outgoing-radiation-refinement}
\cite{stratton1941}. Dual pairings on sources and fields are
Eqs.~\eqref{eq:ch2.dual-source-pairing} and~\eqref{eq:ch2.maxwell-adjoint-pairing}
\cite{stratton1941}. regime qualification Eq.~\eqref{eq:ch2.regime-operator-screen} governs transport
reports on any recovered radiative component \cite{harrington2001}.
Radiation null space (locked definition). Physically, $\ker\mathcal{R}_\omega$ is the set of impressed currents that produce no qualified
far-zone pattern: the wattmeter reading is null even when near probes are not
\cite{stratton1941}. Formally, the \emph{radiation null space} is the source-side kernel
\begin{equation}
\label{eq:ch2.radiation-nullspace-def}
\ker\mathcal{R}_\omega
=
\left\{
\Jimp\in\mathcal{J}_\Omega:
\mathcal{R}_\omega[\Jimp]=\mathbf{0}\in\mathcal{F}_{\mathrm{rad}},
\ \mathcal{D}_{J}[\Jimp]=1
\right\},
\end{equation}
equivalently \(\Pi_{\mathrm{rad}}\,\mathcal{S}_\omega[\Jimp]=\mathbf{0}\) on the outgoing
branch \cite{stratton1941}. Equation~\eqref{eq:ch2.radiation-nullspace-def}
is the Subsection~\ref{sec:ch271} development of Eq.
Eq.~\eqref{eq:ch2.ker-R-omega}: $\mathcal{D}_J=1$ is explicit, and membership is
defined on \(\mathcal{J}_\Omega\) rather than on formal test directions
\cite{stratton1941}. Elements of \(\ker\mathcal{R}_\omega\) produce no qualified far-zone
export through Eq.~\eqref{eq:ch2.radiation-integral-outgoing-upgrade} yet may satisfy
\(\mathcal{S}_\omega[\Jimp]\neq\mathbf{0}\) Eq.~\eqref{eq:ch2.nr-image-of-null-space}
\cite{harrington2001}. Consequence for inversion: any far-zone datum fixes
\(\mathcal{R}_\omega[\Jimp]\) only up to additions in \(\ker\mathcal{R}_\omega\). Linearity
Eq.~\eqref{eq:ch2.radiation-linearity} makes \(\ker\mathcal{R}_\omega\) a linear subspace of
\(\mathcal{J}_\Omega\) closed under admissible superposition \cite{stratton1941}.
Adjoint null space and cokernel. Targets in \(\mathrm{coker}\,\mathcal{R}_\omega\) are far-zone patterns orthogonal to every
reachable image; they obstruct solvability even when no near-field anomaly is visible
\cite{stratton1941}. Let \(\mathcal{R}_\omega^{\dagger}:\mathcal{F}_{\mathrm{rad}}\to\mathcal{J}_\Omega\) denote
the adjoint induced by Eqs.~\eqref{eq:ch2.dual-source-pairing} and
Eq.~\eqref{eq:ch2.maxwell-adjoint-pairing} defines the adjoint pairing. Define
\begin{equation}
\label{eq:ch2.radiation-adjoint-nullspace}
\ker\mathcal{R}_\omega^{\dagger}
=
\left\{
\mathbf{w}\in\mathcal{F}_{\mathrm{rad}}:
\langle \mathbf{w},\,\mathcal{R}_\omega[\Jimp]\rangle_{\mathcal{U}}=0
\ \ \forall\ \Jimp\in\mathcal{J}_\Omega
\right\},
\end{equation}
and the cokernel quotient
\begin{equation}
\label{eq:ch2.radiation-cokernel-def}
\mathrm{coker}\,\mathcal{R}_\omega
=
\mathcal{F}_{\mathrm{rad}}\big/\operatorname{ran}\mathcal{R}_\omega,
\end{equation}
with \(\operatorname{ran}\mathcal{R}_\omega\subset\mathcal{F}_{\mathrm{rad}}\) the radiative
image \cite{stratton1941}. Equations~\eqref{eq:ch2.radiation-adjoint-nullspace}--\eqref{eq:ch2.radiation-cokernel-def}
pair with \(\ker\mathcal{R}_\omega\) through reciprocity-based solvability
Eq.~\eqref{eq:ch2.radiation-solvability-criterion} imported from Subsection~\ref{sec:ch263}:
targets \(\mathbf{g}\in\mathcal{F}_{\mathrm{rad}}\) orthogonal to
\(\ker\mathcal{R}_\omega^{\dagger}\) are reachable, while obstructions lie in cokernel
directions \cite{stratton1941}. Biorthogonality Eq.~\eqref{eq:ch2.radiation-biorthogonality}
expresses the same pairing on modal coordinates without re-deriving adjoint modes here
\cite{harrington2001}.
When \(\mathcal{R}_\omega\) is Fredholm in the sense of Fredholm structure of radiation operators (Subsection~\ref{sec:ch263}),
\begin{equation}
\label{eq:ch2.fredholm-nullspace-alternative}
\dim\ker\mathcal{R}_\omega<\infty
\ \Longleftrightarrow\
\dim\ker\mathcal{R}_\omega^{\dagger}<\infty,
\end{equation}
and index Eq.~\eqref{eq:ch2.radiation-fredholm-index-def} counts the imbalance
\cite{stratton1941}. Equation~\eqref{eq:ch2.fredholm-nullspace-alternative}
does not re-prove Fredholm structure of radiation operators (Subsection~\ref{sec:ch263}); it records the
Chapter~\ref{ch:02} null-space bookkeeping used in inversion
Eq.~\eqref{eq:ch2.inverse-source-equation} \cite{stratton1941}. On exterior domains with
compact \(\Omega_s\), \(\dim\ker\mathcal{R}_\omega\) is typically infinite even when
\(\mathcal{R}_\omega\) is Fredholm on truncated subspaces
Eq.~\eqref{eq:ch2.spectral-compression-operator} \cite{harrington2001}.
On the orthonormal basis of spectral decomposition of the radiation operator (Subsection~\ref{sec:ch261}), define
the \emph{modal null index set}
\begin{equation}
\label{eq:ch2.modal-null-index-set}
\mathcal{N}_{\mathrm{null}}
=
\left\{
n\in\mathcal{N}_{\mathrm{rad}}:
\lambda_n=0
\right\},
\end{equation}
with \(\lambda_n\) from Eq.~\eqref{eq:ch2.radiation-eigenpair-def}
\cite{stratton1941}. Equation~\eqref{eq:ch2.modal-null-index-set} is the
spectral slice of \(\ker\mathcal{R}_\omega\) on declared eigenmodes:
\begin{equation}
\label{eq:ch2.modal-null-sector-reduction}
\sum_{n\in\mathcal{N}_{\mathrm{null}}}
c_n[\mathfrak{D}]\,\Jimp^{(n)}
\in\ker\mathcal{R}_\omega
\quad\text{when admissible superpositions exist},
\end{equation}
using source coefficients Eq.~\eqref{eq:ch2.source-modal-amplitude} and eigenpairs
Eq.~\eqref{eq:ch2.radiation-eigenpair-def} \cite{stratton1941}. Equation
\eqref{eq:ch2.modal-null-sector-reduction} does not exhaust \(\ker\mathcal{R}_\omega\) on
exterior problems because \(\mathcal{N}_{\mathrm{rad}}\) need not span all invisible source
degrees \cite{harrington2001}. modal invertibility on the radiative range (Subsection~\ref{sec:ch263})
already states that coefficients on \(\ker\mathcal{R}_\omega\) remain free in
Eq.~\eqref{eq:ch2.modal-particular-solution}; Subsection~\ref{sec:ch271} identifies the
modal-labeled portion \cite{stratton1941}.
\begin{figure}[t]
\centering
\ChOneFigBlock{%
\begin{tikzpicture}[ch1fig, node distance=5mm and 7mm]
 \node[ch1box] (j) {$\mathcal{J}_\Omega$};
 \node[ch1box, right=7mm of j] (r) {$\mathcal{R}_\omega$};
 \node[ch1box, right=7mm of r] (f) {$\mathcal{F}_{\mathrm{rad}}$};
 \node[ch1box, below=5mm of j] (ker) {$\ker\mathcal{R}_\omega$};
 \node[ch1box, below=5mm of f] (adj) {$\ker\mathcal{R}_\omega^{\dagger}$};
 \draw[ch1flow] (j) -- (r) -- (f);
 \draw[ch1flow, dashed] (ker.south) -- (j.north);
 \draw[ch1flow, dashed] (adj.south) -- (f.north);
\end{tikzpicture}%
}
\caption{Subsection~\ref{sec:ch271} null-space measurement sequence: source kernel
Eq.~\eqref{eq:ch2.radiation-nullspace-def}, radiative image, and adjoint kernel
Eq.~\eqref{eq:ch2.radiation-adjoint-nullspace} paired by solvability
Eq.~\eqref{eq:ch2.radiation-solvability-criterion}.}
\label{fig:ch2.radiation-nullspace-measurement chain}
\end{figure}
Figure interpretation. Read $\mathcal{J}_\Omega$ as the impressed-source Hilbert space and $\mathcal{F}_{\mathrm{rad}}$ as the outgoing radiative image. The forward arrow $\mathcal{R}_\omega=\Pi_{\mathrm{rad}}\circ\mathcal{S}_\omega$ is the only laboratory-export channel; dashed $\ker\mathcal{R}_\omega$ feeds back into $\mathcal{J}_\Omega$ as invisible source degrees that annihilate far-zone patterns Eq.~\eqref{eq:ch2.angular-spectrum-nullspace-law}. The lower adjoint kernel $\ker\mathcal{R}_\omega^{\dagger}$ pairs with cokernel directions through reciprocity-based solvability Eq.~\eqref{eq:ch2.radiation-solvability-criterion}: a nonzero cokernel obstruction means no impressed current can realize the target far field, independent of mesh resolution \cite{stratton1941}. Figure~\ref{fig:ch2.radiation-nullspace-measurement chain} separates kernel and cokernel data before observation composition Eq.~\eqref{eq:ch2.observation-fredholm-composition}.
\phantomsection\label{thm:ch2.radiation-nullspace-structure}
Let \(\mathcal{R}_\omega\) satisfy Eq.~\eqref{eq:ch2.radiation-nullspace-def} on admissible
\(\mathcal{J}_\Omega\) with outgoing radiation condition. Then:
\begin{enumerate}
\item[(i)] \(\ker\mathcal{R}_\omega\) is a linear subspace of \(\mathcal{J}_\Omega\) and
\(\mathcal{R}_\omega[\Jimp]=\mathbf{0}\) iff \(\Jimp\in\ker\mathcal{R}_\omega\).
\item[(ii)] Solvability of \(\mathcal{R}_\omega[\Jimp]=\mathbf{g}\) is equivalent to
Eq.~\eqref{eq:ch2.radiation-solvability-criterion} on declared pairings.
\item[(iii)] On the spectral basis of Subsection~\ref{sec:ch261},
\(\mathcal{N}_{\mathrm{null}}\) Eq.~\eqref{eq:ch2.modal-null-index-set} labels modes with
\(P_n\mathcal{R}_\omega=\mathbf{0}\); modal superpositions
Eq.~\eqref{eq:ch2.modal-null-sector-reduction} lie in \(\ker\mathcal{R}_\omega\) when
admissible.
\item[(iv)] Domain class fixes dimensionality:
\begin{equation}
\label{eq:ch2.nullspace-dimensional-classification}
\begin{aligned}
\text{bounded cavity + discrete book:}\ &
\dim\ker\mathcal{R}_\omega\ \text{finite on truncated}\ \mathcal{R}_\omega^{(N)};\\
\text{exterior + compact}\ \Omega_s:\ &
\dim\ker\mathcal{R}_\omega=\infty\ \text{typ. on full}\ \mathcal{J}_\Omega.
\end{aligned}
\end{equation}
\end{enumerate}. Part (i) follows from linearity Eq.~\eqref{eq:ch2.radiation-linearity} and definition
Eq.~\eqref{eq:ch2.radiation-nullspace-def} \cite{stratton1941}. Part (ii) is Fredholm structure of radiation operators (Subsection~\ref{sec:ch263}) and
Eq.~\eqref{eq:ch2.radiation-solvability-criterion}, cited without re-proof
\cite{harrington2001}. Part (iii) uses Eq.~\eqref{eq:ch2.modal-fredholm-reduction} and
\(\lambda_n=0\) in Eq.~\eqref{eq:ch2.radiation-eigenpair-def} \cite{stratton1941}. Part (iv)
contrasts discrete cavity spectra Eq.~\eqref{eq:ch2.cavity-discrete-radiation-spectrum} with
continuous exterior angular content Eq.~\eqref{eq:ch2.exterior-continuous-spectral-integral}
\cite{jackson1999}.
\phantomsection\label{prop:ch2.null-kernel-field-image}
If \(\Jimp\in\ker\mathcal{R}_\omega\), then radiative and non-radiative components satisfy
\(\mathbf{u}_{\mathrm{rad}}=\mathbf{0}\) and
\(\mathbf{u}_{\mathrm{nr}}=\mathcal{S}_\omega[\Jimp]\) from
Eqs.~\eqref{eq:ch2.radiative-component-def}--\eqref{eq:ch2.nonradiative-component-def};
near-zone energy in \(\mathbf{u}_{\mathrm{nr}}\) may be large while qualified export vanishes
\cite{stratton1941}.
Composition with spectral and geometric filters. Null-space membership implies modal annihilation on exporting sectors:
\phantomsection\label{eq:ch2.geometric-modal-annihilation}
\begin{equation}
\label{eq:ch2.nullspace-modal-annihilation-law}
\Jimp\in\ker\mathcal{R}_\omega
\ \Longrightarrow\
a_n(\Jimp)=0
\ \text{for all}\ n\ \text{with}\ \lambda_n\neq 0,
\end{equation}
with far coefficients \(a_n\) from Eq.~\eqref{eq:ch2.spectral-far-pattern-expansion}
\cite{stratton1941}. Equation~\eqref{eq:ch2.nullspace-modal-annihilation-law}
expresses Eq.~\eqref{eq:ch2.nullspace-spectral-composition} from Section~\ref{sec:ch27} without
re-deriving the modal expansion \cite{stratton1941}. Geometric constraints
geometric constraints on radiative spectra (Subsection~\ref{sec:ch264}) may force additional \(a_n=0\) for sources
outside \(\ker\mathcal{R}_\omega\) \cite{harrington2001}. Composition law
Eq.~\eqref{eq:ch2.constraint-composition-law} budgets identifiability deficits separately
\cite{stratton1941}.
Fredholm index Eq.~\eqref{eq:ch2.radiation-fredholm-index-def} is unchanged under compact
perturbations Eq.~\eqref{eq:ch2.radiation-fredholm-decomposition} and under Tikhonov
regularization Eq.~\eqref{eq:ch2.tikhonov-source-reconstruction} on the range side
\cite{stratton1941}. Regularization stabilizes
Eq.~\eqref{eq:ch2.modal-particular-solution} for small \(|\lambda_n|\) but leaves
\(\ker\mathcal{R}_\omega\) unchanged: invisible source components cannot be recovered from
far-zone data at any regularization level \cite{stratton1941}. Compression truncation
Eq.~\eqref{eq:ch2.compression-error-bound} approximates \(\operatorname{ran}\mathcal{R}_\omega\),
not \(\ker\mathcal{R}_\omega\) \cite{harrington2001}.
Null-space structure of radiation operators (Subsection~\ref{sec:ch271}) assumes linear media at fixed \(\omega\),
outgoing radiation condition, and admissible \(\mathcal{J}_\Omega\) Eq.~\eqref{eq:ch2.source-admissibility}
\cite{stratton1941}. Loss modifies \(\lambda_n\) and may merge modal null sectors
without promoting evanescent content to export through
Eq.~\eqref{eq:ch2.rad-null-on-evanescent} \cite{stratton1941}. Nonlinear or time-varying
sources require the causal chain of Subsection~\ref{sec:ch223}; harmonic null spaces here are
fixed-frequency slices \cite{harrington2001}. regime qualification
Eq.~\eqref{eq:ch2.regime-component-criterion} remains mandatory before interpreting near-zone
reactive dominance as null-space membership \cite{stratton1941}.
\medskip
\noindent
Subsection~\ref{sec:ch271} fixed operator null-space structure through
Eq.~\eqref{eq:ch2.radiation-nullspace-def}, adjoint and cokernel pairings, and the modal null
sector Eqs.~\eqref{eq:ch2.modal-null-index-set}--\eqref{eq:ch2.modal-null-sector-reduction}
\cite{stratton1941}. Constructive non-radiating current families that realize the kernel in the
laboratory develop next.
\phantomsection\label{sec:ch272}
